\documentclass[11pt,a4paper]{article}
\usepackage{../../../_templates/latex/mastery-style}

\newcommand{\QID}[1]{\textcolor{MPrimary}{\textbf{[#1]}}}
\newcommand{\Tag}[1]{\textcolor{MMuted}{\texttt{#1}}}

\begin{document}

\MasterySetBrand{25Maths}{Mastery Series}
\MasterySetHeader{Quadratics Mastery Pack}{Module 2: Full Solutions}
\MasterySetFooter{Quadratics v1 | Module 2 | Solutions}

\MasteryModuleCover
  {2}
  {Full Solutions}
  {Provide step-by-step reasoning and method checkpoints for every workbook question}

\section{Marking Standard}
\begin{enumerate}[leftmargin=1.3em]
  \item Method choice is visible.
  \item Algebra steps are logically ordered.
  \item All roots are listed on final line.
  \item Context constraints are checked where needed.
\end{enumerate}

\begin{keyconceptbox}
Treat each solution as a reusable template. Students should copy both the method choice and final-answer format.
\end{keyconceptbox}

\section{Block A Solutions (Q01-Q08)}

\subsection*{\QID{Q01} Solve \(x^2-9=0\)}
\textbf{Method:} difference of squares.
\[
x^2-9=0 \Rightarrow (x-3)(x+3)=0
\]
\[
\therefore x=3\;\text{or}\;x=-3
\]
\textbf{Final answer:} \(x=\pm3\). \Tag{QP-01-missing-root}

\subsection*{\QID{Q02} Solve \(x^2+7x+12=0\)}
\textbf{Method:} factorising.
\[
x^2+7x+12=(x+3)(x+4)=0
\]
\[
\therefore x=-3\;\text{or}\;x=-4
\]
\textbf{Final answer:} \(x=-3,-4\). \Tag{QP-01-missing-root}

\subsection*{\QID{Q03} Solve \(x^2-11x+24=0\)}
\textbf{Method:} factorising.
\[
x^2-11x+24=(x-3)(x-8)=0
\]
\[
\therefore x=3\;\text{or}\;x=8
\]
\textbf{Final answer:} \(x=3,8\). \Tag{QP-01-missing-root}

\subsection*{\QID{Q04} Solve \(2x^2-18=0\)}
\textbf{Method:} simplify then solve.
\[
2x^2=18 \Rightarrow x^2=9 \Rightarrow x=\pm3
\]
\textbf{Final answer:} \(x=\pm3\). \Tag{QP-02-sign-slip}

\subsection*{\QID{Q05} Solve \(x^2+2x-15=0\)}
\textbf{Method:} factorising.
\[
x^2+2x-15=(x+5)(x-3)=0
\]
\[
\therefore x=-5\;\text{or}\;x=3
\]
\textbf{Final answer:} \(x=-5,3\). \Tag{QP-05-arithmetic-load}

\subsection*{\QID{Q06} Solve \(x^2=5x\)}
\textbf{Method:} rearrange then factorise.
\[
x^2-5x=0 \Rightarrow x(x-5)=0
\]
\[
\therefore x=0\;\text{or}\;x=5
\]
\textbf{Final answer:} \(x=0,5\). \Tag{QP-02-sign-slip}

\subsection*{\QID{Q07} Solve \((x-4)(x+1)=0\)}
\textbf{Method:} zero product rule.
\[
x-4=0\Rightarrow x=4,\quad x+1=0\Rightarrow x=-1
\]
\textbf{Final answer:} \(x=4,-1\). \Tag{QP-01-missing-root}

\subsection*{\QID{Q08} Solve \(x^2-4x=45\)}
\textbf{Method:} rearrange and factorise.
\[
x^2-4x-45=0
\]
\[
(x-9)(x+5)=0
\]
\[
\therefore x=9\;\text{or}\;x=-5
\]
\textbf{Final answer:} \(x=9,-5\). \Tag{QP-02-sign-slip}

\section{Block B Solutions (Q09-Q16)}

\subsection*{\QID{Q09} Solve \(x^2-6x+1=0\)}
\textbf{Method:} formula.
\[
x=\frac{-(-6)\pm\sqrt{(-6)^2-4(1)(1)}}{2(1)}=\frac{6\pm\sqrt{32}}{2}=3\pm2\sqrt2
\]
\textbf{Final answer:} \(x=3\pm2\sqrt2\). \Tag{QP-03-method-mismatch}

\subsection*{\QID{Q10} Solve \(3x^2+2x-5=0\)}
\textbf{Method:} formula then simplify.
\[
x=\frac{-2\pm\sqrt{2^2-4(3)(-5)}}{2\cdot3}=\frac{-2\pm\sqrt{64}}{6}=\frac{-2\pm8}{6}
\]
\[
\therefore x=1\;\text{or}\;x=-\frac{5}{3}
\]
\textbf{Final answer:} \(x=1,-\frac{5}{3}\). \Tag{QP-03-method-mismatch}

\subsection*{\QID{Q11} Solve \(5x^2-3x-2=0\)}
\textbf{Method:} formula.
\[
x=\frac{3\pm\sqrt{(-3)^2-4(5)(-2)}}{2\cdot5}=\frac{3\pm\sqrt{49}}{10}=\frac{3\pm7}{10}
\]
\[
\therefore x=1\;\text{or}\;x=-\frac{2}{5}
\]
\textbf{Final answer:} \(x=1,-\frac25\). \Tag{QP-05-arithmetic-load}

\subsection*{\QID{Q12} Solve \(2x^2-4x+1=0\), 3 s.f.}
\textbf{Method:} formula and rounding.
\[
x=\frac{4\pm\sqrt{(-4)^2-4(2)(1)}}{2\cdot2}=\frac{4\pm\sqrt8}{4}=1\pm\frac{\sqrt2}{2}
\]
\[
x\approx1.707\;\text{or}\;0.293
\]
\textbf{Final answer:} \(x\approx1.71,0.293\) (3 s.f.). \Tag{QP-05-arithmetic-load}

\subsection*{\QID{Q13} Solve \(x^2+4x-1=0\)}
\textbf{Method:} complete square.
\[
x^2+4x=1
\]
\[
(x+2)^2-4=1 \Rightarrow (x+2)^2=5
\]
\[
x+2=\pm\sqrt5 \Rightarrow x=-2\pm\sqrt5
\]
\textbf{Final answer:} \(x=-2\pm\sqrt5\). \Tag{QP-03-method-mismatch}

\subsection*{\QID{Q14} Solve \(x^2+10x+20=0\)}
\textbf{Method:} complete square.
\[
x^2+10x=-20
\]
\[
(x+5)^2-25=-20 \Rightarrow (x+5)^2=5
\]
\[
x+5=\pm\sqrt5 \Rightarrow x=-5\pm\sqrt5
\]
\textbf{Final answer:} \(x=-5\pm\sqrt5\). \Tag{QP-05-arithmetic-load}

\subsection*{\QID{Q15} Show solutions to \(x^2-4x-1=0\) are \(2\pm\sqrt5\)}
\textbf{Method:} complete square proof.
\[
x^2-4x=1
\]
\[
(x-2)^2-4=1 \Rightarrow (x-2)^2=5
\]
\[
x-2=\pm\sqrt5 \Rightarrow x=2\pm\sqrt5
\]
\textbf{Final answer:} shown. \Tag{QP-03-method-mismatch}

\subsection*{\QID{Q16} Determine number of real roots of \(4x^2+4x+5=0\)}
\textbf{Method:} discriminant test.
\[
\Delta=b^2-4ac=4^2-4(4)(5)=16-80=-64<0
\]
\textbf{Final answer:} no real roots. \Tag{QP-06-discriminant-misread}

\section{Block C Solutions (Q17-Q24)}

\subsection*{\QID{Q17} Rectangle area model \(x(x+3)=54\)}
\textbf{Method:} model, solve, context filter.
\[
x^2+3x-54=0
\]
\[
(x+9)(x-6)=0 \Rightarrow x=-9\;\text{or}\;x=6
\]
Length must be positive, so \(x=6\).
\textbf{Final answer:} \(x=6\,\text{cm}\). \Tag{QP-04-invalid-context-root}

\subsection*{\QID{Q18} \(h=t(12-3t)\), find when \(h=0\)}
\textbf{Method:} zero product rule.
\[
t(12-3t)=0
\]
\[
t=0\;\text{or}\;12-3t=0\Rightarrow t=4
\]
\textbf{Final answer:} \(t=0\) and \(t=4\). \Tag{QP-04-invalid-context-root}

\subsection*{\QID{Q19} Roots \(3\) and \(-2\), form equation}
\textbf{Method:} root factors.
\[
(x-3)(x+2)=0
\]
\[
x^2-x-6=0
\]
\textbf{Final answer:} \(x^2-x-6=0\). \Tag{QP-01-missing-root}

\subsection*{\QID{Q20} Monic quadratic with repeated root \(x=3\)}
\textbf{Method:} repeated root form.
\[
(x-3)^2=0
\]
\[
x^2-6x+9=0
\]
\textbf{Final answer:} \((x-3)^2=0\) or \(x^2-6x+9=0\). \Tag{QP-03-method-mismatch}

\subsection*{\QID{Q21} For \(y=x^2-6x+5\), find axis, vertex, and intercepts}
\textbf{Method:} complete square and factor.
\[
y=x^2-6x+5=(x-3)^2-4
\]
Axis of symmetry: \(x=3\). Vertex: \((3,-4)\).
\[
x^2-6x+5=0 \Rightarrow (x-1)(x-5)=0
\]
So \(x\)-intercepts are \((1,0)\) and \((5,0)\).
\textbf{Final answer:} axis \(x=3\), vertex \((3,-4)\), intercepts \((1,0),(5,0)\). \Tag{QP-06-discriminant-misread}

\subsection*{\QID{Q22} \(x^2-kx+12=0\), one root is 3}
\textbf{Method:} product and sum of roots.
\[
\text{If roots are }3\text{ and }r,\quad 3r=12\Rightarrow r=4
\]
\[
\text{Sum of roots}=3+4=7=k
\]
\textbf{Final answer:} other root \(4\), and \(k=7\). \Tag{QP-06-discriminant-misread}

\subsection*{\QID{Q23} Solve \(x^2-5x+6\le0\)}
\textbf{Method:} factor and sign interval.
\[
x^2-5x+6=(x-2)(x-3)
\]
Upward-opening quadratic is non-positive between roots.
\textbf{Final answer:} \(2\le x\le3\).\\
\Tag{QP-03-method-mismatch}

\subsection*{\QID{Q24} Consecutive integers product 156}
\textbf{Method:} model and context filter.
Let integers be \(x\) and \(x+1\).
\[
x(x+1)=156 \Rightarrow x^2+x-156=0
\]
\[
(x+13)(x-12)=0 \Rightarrow x=-13\;\text{or}\;x=12
\]
Positive pair is \(12\) and \(13\).
\textbf{Final answer:} \(12,13\). \Tag{QP-04-invalid-context-root}

\begin{tcolorbox}[masterycard,colback=MSoft]
\textbf{Marking handoff:} export repeated error tags to Module 3 and assign targeted drills before any new mixed test.
\end{tcolorbox}

\end{document}
